\section{Introduction}
Vision--language models support general-purpose visual understanding, while decision pipelines require low-latency judgments over runtime-defined alternatives. The System One paradigm addresses this requirement through direct probabilistic prediction, extending naturally to image-conditioned decisions~\citep{yu_2026_visualjev,zhou_2026_pixeljev}.

Candidate scoring must accommodate both order symmetry and set dependence: reordering should preserve candidate identity, while changing the alternatives may change relative preferences. Shared encoding and joint comparison combine these properties through set representations~\citep{zaheer_2017_deepsets,lee_2019_set}.

We present OpenVisualJev\footnote{An independently developed research implementation of a Jev-style multimodal decision model.}, which separates candidate grounding from joint comparison. Tree-masked attention shares image and question prefixes across isolated candidate branches in one forward pass. A position-free set module compares their representations, and a shared pointer produces candidate probabilities. Calibration-aware training and held-out temperature scaling complete the model (Figure~\ref{fig:set-pointer-method}).

Controlled order and relation tests examine the architecture's structural properties; unified recognition and natural VQA assess task performance. Ablations isolate candidate interaction, training with variable candidate counts, and calibration.

Our contributions are threefold:
\begin{itemize}
\item A multimodal decision model combining pretrained visual representations, dynamic candidate scoring, and probability calibration.
\item A permutation-equivariant scoring architecture with joint candidate comparison, together with an analysis of its structural properties.
\item Controlled structural tests, recognition and VQA evaluations, and ablations of interaction, candidate-count training, and calibration.
\end{itemize}
